%=======================================================================================
% METHODOLOGY REDRAFT -- 2026-08-19
% This is a separate copy. The original draft is unchanged.
% Numerical findings belong in Results and Discussion, not here.
% Long feature lists, search grids and secondary checks belong in the Appendix.
%=======================================================================================

\chapter{Methodology and experimental design}
\label{ch:methodology}

This chapter describes how the data in Chapter~\ref{ch:data} were used to answer the research
questions. It defines the splits, feature sets, models, targets and evaluation procedures. The
same held-out observations were used when models were compared.

\section{Experimental framework}
\label{sec:method-framework}

RQ1 tests predictive accuracy. RQ2a tests when environmental conditioning corrects departure
from a shared growth curve. RQ2b attributes persistent departure from that curve. RQ3 attributes
plot-specific asymptote deviation and tests whether the associations vary over space.

\begin{table}[!htbp]
\centering
\scriptsize
\renewcommand{\arraystretch}{1.15}
\begin{tabular}{p{0.9cm}p{3.5cm}p{3.2cm}p{2.8cm}}
\toprule
RQ & Unit and target & Models or comparison & Main evaluation \\
\midrule
RQ1 & Plot--survey row; 95th-percentile top height & CR, linear regression, random forest,
XGBoost, DNN, PINN and PINN-$k$ & Five-fold spatial block CV; temporal and plot-level checks \\
RQ2a & Plot--survey row; error reduction from the fold-matched CR curve & Conditioned
DNN, PINN, PINN-$k$ and XGBoost; no-env and permutation controls & Five-fold spatial block
CV; CR-error quartiles \\
RQ2b & Plot; mean pooled-CR residual & Elastic Net, XGBoost and two-stage
mixed-effects model & Five-fold spatial block CV; coefficients, SHAP and residual Moran's $I$ \\
RQ3 & Plot; difference between fitted and yield-class-implied asymptotes & Elastic Net, XGBoost
and GNNWR & Five-fold spatial block CV; attribution, Moran's $I$ and target audit \\
\bottomrule
\end{tabular}
\caption{Units, targets and evaluation used for each research question.}
\label{tab:rq-framework}
\end{table}

% Keep one high-level pipeline in the main text. Move any detailed pipeline to the Appendix.
\begin{figure}[!htbp]
\centering
\includegraphics[width=\textwidth]{figures/fig_methodology/method_fig_pipeline-overview.png}
\caption{Analysis pipeline. The two cohorts feed shared spatial partitions and separate
RQ-specific feature sets. RQ1 predicts top height, RQ2 examines departure from a shared curve,
and RQ3 examines plot-specific asymptote deviation.}
\label{fig:pipeline-overview}
\end{figure}

\section{Data partitions and leakage control}
\label{sec:method-splits}

The four-survey cohort is the main cohort. The six-survey cohort is a nested sensitivity cohort
with longer plot histories but fewer plots and compartments. RQ1 uses plot--survey rows. RQ2b and
RQ3 use one row per plot. RQ2a evaluates RQ1 predictions at plot--survey level.

\subsection{Evaluation splits}

The spatial block design is the main evaluation. Whole compartments are assigned to folds, so a
compartment cannot occur in both training and test data. A 60\,m buffer provides a second guard at
held-out boundaries: training plots within 60\,m of validation or test plots are removed, while
held-out plots are retained \citep{roberts2017Cross_C5_EVAL}.

For five-fold spatial cross-validation, compartments are balanced by row count. One fold is used
for testing, the next for validation and the other three for training. Each compartment is tested
once. Held-out predictions from all folds are pooled for population-level metrics
\citep{ploton2020Spatial_C5_EVAL}.

Two supporting RQ1 splits were retained. The plot-level split assigns plots at random in a
60/20/20 ratio and gives an easier, spatially leaky reference. The temporal split trains on 2008
and 2012, validates on 2021 and tests on 2023 for the four-survey cohort. For the six-survey cohort,
2002--2012 are used for training, 2021 for validation and 2023 for testing
\citep{lynch2025Digital_C4_NN}. All primary split assignments use seed 42.

\subsection{Preprocessing}

Continuous variables are standardised using the training partition only. The fitted transform is
then applied unchanged to validation and test data. Age is scaled separately because its training
standard deviation is needed to return the PINN derivative to metres per year. Categorical soil
classes are one-hot encoded. If a required variable is missing, the whole plot trajectory is
removed rather than retaining an incomplete history. Missing-value counts and plot identifiers
are reported in Appendix~\ref{app:missing-cells}.

Compartment, sub-compartment and block identifiers are excluded from all model inputs. Two
neighbour-height variables were also excluded because they were calculated from nearby plots'
observed heights before splitting and therefore leaked held-out outcomes.

\section{Environmental feature sets}
\label{sec:feature-sets}

Feature sets are constructed separately for RQ1, RQ2b and RQ3 because their targets differ. The
same four-step rule is used in each case.

\begin{enumerate}
\setlength\itemsep{0pt}
\item Variables with insufficient coverage are removed before ranking.
\item Deterministic duplicates are removed. For pairs with Spearman $|\rho|\geq0.95$, one
representative is retained, favouring externally measured variables over derived variables.
\item Candidates are ranked by the mean rank from univariate Spearman $|\rho|$, permutation
importance and drop-column change in held-out $R^2$.
\item Variance inflation factors are reduced iteratively to a threshold of 5
\citep{obrien2007Caution_C5_EVAL}. Protected stand variables remain in the matrix but still
contribute to the VIF of other variables.
\end{enumerate}

The protected baseline comprises canopy cover, time since thinning, its missingness indicator and
recent thinning. Set~1 contains these variables only. Set~2 adds the ten highest-ranked candidates.
Set~3 adds the upper half of the terrain and wind candidates. Set~4 extends Set~3 with the upper
half of the climate, soil and spatial-context candidates. Exact membership and VIF diagnostics are
given in Appendix~\ref{app:feature-sets}.

Feature membership is fixed once from the designated training partition rather than rebuilt in
each fold. This keeps the model comparisons on the same inputs. It also means the reported fold
variation measures model and spatial-split variation, not variation in feature selection.

\section{RQ1: Top-height prediction}
\label{sec:rq1-method}

RQ1 compares a biological curve, standard tabular models and neural models. The headline
comparison uses Set~3 and the same five spatial folds. The other feature sets test whether adding
environmental inputs changes the result.

\subsection{Chapman--Richards baseline}

The Chapman--Richards (CR) curve is
\begin{equation}
H(a)=y_{\max}\left(1-e^{-ka}\right)^p,
\label{eq:cr}
\end{equation}
where $H(a)$ is height at age $a$, $y_{\max}$ is the asymptote, $k$ is the growth rate and $p$
controls curve shape. Non-linear least squares is run from five starting points. The fit with the
lowest sum of squared errors is retained.

The CR curve is refitted on the training rows of every split and fold. Its parameters are then
frozen. Validation and test observations therefore do not affect the baseline or the PINN
guidance derived from it.

\subsection{Conventional models}

Linear regression provides a low-complexity baseline. Random forest provides a bagged tree
ensemble. XGBoost provides a boosted tree ensemble \citep{chen2016XGBoost_C4_NN}. These models use
the same RQ1 feature set and held-out folds.

RQ1 XGBoost hyperparameters were selected for the top-height target. Twenty-seven candidate
configurations were compared across five folds and both cohorts, giving 270 fits. Selection used
validation $R^2$ only; test results were read after selection. The search space and selected values
are given in Appendix~\ref{app:hyperparameters}. RQ2b and RQ3 use the existing fixed XGBoost
configuration instead: depth 4, learning rate 0.04 and at most 500 boosting rounds with early
stopping. Those attribution models were not retuned.

\subsection{DNN and physics-informed models}
\label{sec:rq1-neural}

All neural models share a three-layer prediction trunk with 128 units per layer and Leaky-ReLU
activation. Optimisation uses Adam with learning rate $10^{-4}$, weight decay $10^{-5}$, gradient
clipping at 1.0, a maximum of 500 epochs, batch size 256 and early stopping with patience 40. The
DNN minimises
\begin{equation}
\mathcal{L}_{\mathrm{data}}=\frac{1}{N}\sum_{i=1}^{N}
\left(\hat H_i-H_i\right)^2.
\label{eq:data-loss}
\end{equation}

The environment-conditioned DNN concatenates the environmental variables to the main input. The
PINN variants use the same prediction trunk, but environmental variables enter small auxiliary
networks. The architectures therefore share the same core capacity, while the environmental
information is routed differently.

The analytical CR derivative is
\begin{equation}
H'(a)=y_{\max}pk e^{-ka}\left(1-e^{-ka}\right)^{p-1}.
\label{eq:cr-derivative}
\end{equation}
The physics loss compares this derivative with the network derivative obtained by automatic
differentiation:
\begin{equation}
\mathcal{L}_{\mathrm{phys}}=\frac{1}{N}\sum_{i=1}^{N}
\left(\frac{\partial\hat H_i}{\partial a_i}-H'(a_i)\right)^2.
\label{eq:physics-loss}
\end{equation}
The network derivative is divided by the training-set age standard deviation before comparison,
so both terms use metres per year.

The trajectory loss uses consecutive training surveys from the same plot:
\begin{equation}
\mathcal{L}_{\mathrm{traj}}=\frac{1}{M}\sum_{j=1}^{M}
\left(\frac{\hat H_{j,2}-\hat H_{j,1}}{\Delta a_j}-H'(\bar a_j)\right)^2,
\label{eq:trajectory-loss}
\end{equation}
where $\bar a_j$ is the midpoint age. Pairs are formed only when both rows belong to the training
partition. The total loss is
\begin{equation}
\mathcal{L}=\mathcal{L}_{\mathrm{data}}+
\lambda_{\mathrm{phys}}\mathcal{L}_{\mathrm{phys}}+
\lambda_{\mathrm{traj}}\mathcal{L}_{\mathrm{traj}}+
10^{-5}\sum_{\theta}|\theta|.
\label{eq:total-loss}
\end{equation}
The default uses $\lambda_{\mathrm{phys}}=\lambda_{\mathrm{traj}}=1$. Joint weights 0, 1 and 2
are compared in the ablation. Weight 0 is the physics-disabled control.

The first conditioned PINN changes the asymptote:
\begin{equation}
y_{\max}^{(i)}=y_{\max}+f_y(\mathbf{x}_i).
\label{eq:ymax-adjust}
\end{equation}
PINN-$k$ also changes the rate while keeping it positive:
\begin{equation}
k^{(i)}=k\exp\!\left(f_k(\mathbf{x}_i)\right).
\label{eq:k-adjust}
\end{equation}
Both auxiliary networks have one 16-unit hidden layer and start near zero. The shape parameter
$p$ remains fixed in every variant.

\begin{figure}[!htbp]
\centering
\includegraphics[width=\textwidth]{figures/fig_methodology/method_fig_nn_model_architecture.png}
\caption{Shared neural prediction trunk and the CR-guided losses. Environmental inputs enter the
DNN directly, but enter the PINN through the $y_{\max}$ and optional $k$ subnetworks. The physics
loss uses the instantaneous derivative; the trajectory loss uses changes between observed survey
pairs.}
\label{fig:model-architecture}
\end{figure}

\section{RQ2a: Environmental correction of curve departure}
\label{sec:rq2a-method}

RQ2a reuses held-out RQ1 predictions. For test row $i$, improvement over the fold-matched CR curve
is
\begin{equation}
d_i=\left|H_i-\hat H_i^{\mathrm{CR}}\right|-
\left|H_i-\hat H_i^{\mathrm{model}}\right|.
\label{eq:residual-reduction}
\end{equation}
A positive value means that the model is closer to the observed height. The analysis reports mean
$d_i$, its percentage of mean CR error, and the percentage of rows with $d_i>0$.

Within each held-out fold, rows are ranked by absolute CR error and divided into four equal groups.
Q1 contains rows where CR fits best; Q4 contains rows where it fits worst. Mean improvement and the
percentage improved are calculated for each quartile. This is repeated for the environment-
conditioned DNN, PINN, PINN-$k$ and XGBoost across Sets~2--4 and both cohorts. The Set~3 DNN is
evaluated over five training seeds. Other RQ2a rows use seed 42; their reported spread is therefore
fold-to-fold, not seed-to-seed variation.

Two controls isolate the source of any improvement. First, XGBoost is evaluated with and without
environmental inputs. Second, the environmental values are shuffled across held-out rows while
the non-environment inputs are kept fixed. The shuffled control preserves each variable's
distribution but breaks its link to the plot. The same error-reduction calculation is then
repeated. % VERIFY from the script: confirm that shuffling was within each held-out fold and that
% the trained model was not refitted. The Results draft currently implies this design.

The spatial pattern of $d_i$ is assessed with Moran's $I$ using the same weights as the residual
analysis in Section~\ref{sec:evaluation}.

\section{RQ2b: Persistent departure from the shared curve}
\label{sec:rq2b-method}

RQ2b uses the four-survey cohort. The six-survey cohort has too few compartments for the main
spatial attribution analysis. For each plot--survey row,
\begin{equation}
r_{it}=H_{it}-H(a_{it}),
\end{equation}
and the plot target is
\begin{equation}
\bar r_i=\frac{1}{T_i}\sum_{t=1}^{T_i}r_{it}.
\label{eq:mean-cr-residual}
\end{equation}
Positive values mark plots that remain above the pooled curve on average. Age is excluded from the
candidate predictors because it already defines the CR expectation.

Elastic Net and XGBoost are fitted under five-fold compartment CV. Elastic Net selects its penalty
and $\ell_1$ mixing ratio by cross-validation within the training partition. XGBoost uses the fixed
configuration stated in Section~\ref{sec:rq1-method}. Pooled out-of-fold $R^2$ measures held-out
attribution.

A two-stage mixed-effects model is also fitted. Stage~1 constructs $\bar r_i$ from the nonlinear
CR curve. Stage~2 fits
\begin{equation}
\bar r_i=\boldsymbol{\beta}^{\top}\mathbf{x}_i+u_{c[i]}+\varepsilon_i,
\label{eq:mixed-effects}
\end{equation}
where $u_{c[i]}\sim\mathcal{N}(0,\sigma_u^2)$ is a random intercept for compartment. Random slopes
are not used. Fixed-effect coefficients are estimated by restricted maximum likelihood within
each training fold. Spatial variance explained is the proportional reduction in compartment
random-intercept variance relative to the intercept-only model. Coefficients are reported as the
mean $\pm$ sample SD across folds.

Four checks support interpretation. First, Set~4 is refitted without canopy cover. Second, Pearson
correlations are calculated between canopy cover and each target. Third, SHAP values
are recalculated after the canopy-cover removal to test whether another predictor absorbs its
credit. Fourth, the share of raw variance occurring between compartments is calculated for slope
and TOPEX to help interpret disagreement between the mixed model and Elastic Net.

Elastic Net coefficients, mixed-effects coefficients and XGBoost mean absolute SHAP values provide
three attribution views. SHAP is calculated after model selection and does not affect feature-set
membership. Moran's $I$ is calculated from the pooled held-out Elastic Net and XGBoost residuals.

\section{RQ3: Plot-specific asymptote deviation}
\label{sec:rq3-method}

RQ3 estimates one asymptote per plot while holding the Forestry Commission rate and shape
parameters fixed. For plot $i$, define
\begin{equation}
g_{it}=\left(1-e^{-p_{4,i}a_{it}}\right)^{p_{5,i}}.
\end{equation}
The least-squares fit through the origin is
\begin{equation}
\hat y_{\max,i}=\frac{\sum_t H_{it}g_{it}}{\sum_t g_{it}^2}.
\label{eq:plot-ymax-fit}
\end{equation}
The yield-class-implied asymptote is derived from the plot's Forestry Commission parameters:
\begin{equation}
y_{\max,i}^{\mathrm{yldc}}=\mathrm{yldc}_i p_{2,i}+p_{1,i}+2p_{3,i}.
\end{equation}
The attribution target is
\begin{equation}
\Delta y_{\max,i}=\hat y_{\max,i}-y_{\max,i}^{\mathrm{yldc}}.
\label{eq:local-ymax-diff}
\end{equation}
No environmental variable enters this target. Fixing $p_4$ and $p_5$ avoids fitting three curve
parameters from only four or six observations, but places all shape mismatch into the fitted
asymptote.

Elastic Net and XGBoost are fitted to $\Delta y_{\max}$ for Sets~2--4. GNNWR adds spatially varying
coefficients \citep{du2020Geographically_C6_GWR}:
\begin{equation}
\hat y_i=\beta_0(s_i)+\sum_j\beta_j(s_i)x_{ij}.
\label{eq:gnnwr}
\end{equation}
For each focal plot, the input to the weighting network is the MinMax-scaled vector of Euclidean
distances to every training reference plot. The network outputs one multiplicative weight for
each coefficient and the intercept. These weights modify a frozen global OLS fit. The reported
runs use the full, uncapped training reference set. Held-out plots use the training reference set;
no kernel or fixed bandwidth is imposed. All three models use the same five spatial folds.

XGBoost gain is the reduction in training loss produced by splits that use a feature, accumulated
over the ensemble. Mean absolute SHAP measures the size of a feature's contribution to individual
predictions \citep{lundberg2017unified}. Gain describes use during tree construction; SHAP
describes contribution to fitted predictions. RQ3 SHAP values are calculated out of fold and then
pooled.

\subsection{Target and outlier audit}

The curve-fit population excludes a plot if any consecutive survey interval meets either of these
rules:
\begin{itemize}
\setlength\itemsep{0pt}
\item clearfell-like: height falls by at least 25\%, canopy cover by at least 70\%, and volume by
at least 70\%;
\item measurement-inconsistent: height falls by at least 25\%, while canopy cover falls by no more
than 10\% and volume does not fall.
\end{itemize}
Intervals with a height collapse but neither pattern are labelled ambiguous and retained. Exact
counts are reported in Results rather than here.

Extreme fitted asymptotes are flagged separately using Tukey fences,
\begin{equation}
\left[Q_1-1.5\,\mathrm{IQR},\;Q_3+1.5\,\mathrm{IQR}\right],
\label{eq:tukey-fences}
\end{equation}
calculated within each cohort. % CITATION: add the confirmed BibTeX key for Tukey (1977).
The flagged plots remain in the main analysis. A removal sensitivity recalculates pooled metrics
from the saved held-out predictions after excluding them; the models are not refitted.

The audit then checks whether the same plots recur among the worst 1\% of residuals across models
and sets. It compares their within-plot residual range, distance to compartment boundaries,
compartment concentration and wind-shelter variables with the full population. To test a known
survey-boundary concern, the worst-fitting survey year is counted and each flagged plot is refitted
after removing its 2008 observation. These checks diagnose the derived target; they do not prove
an ecological cause.

\section{Evaluation and robustness}
\label{sec:evaluation}

For residual $e_i=H_i-\hat H_i$, the main metrics are
\begin{align}
\mathrm{RMSE}&=\sqrt{\frac{1}{N}\sum_i e_i^2}, &
\mathrm{MAE}&=\frac{1}{N}\sum_i|e_i|, \\
R^2&=1-\frac{\sum_i e_i^2}{\sum_i(H_i-\bar H)^2}, &
\mathrm{Bias}&=\frac{1}{N}\sum_i e_i.
\end{align}
Positive bias denotes under-prediction. Each table states whether it reports a fold mean or a
pooled metric. Fold summaries are the mean $\pm$ sample SD across five folds. Pooled metrics are
calculated once from all out-of-fold predictions. These are not interchangeable.

For the four-survey spatial analyses, 95\% confidence intervals use 2,000 cluster-bootstrap
resamples of whole compartments. The 2.5th and 97.5th percentiles form the interval. This preserves
spatial clustering better than resampling individual rows. The smaller six-survey RQ3 results are
reported as point estimates; GNNWR is also repeated over three split seeds to test sign stability.

A semivariogram measures how quickly residuals become less similar as distance between plots increases; its range is the distance beyond which they are no longer spatially correlated.
Global Moran's $I$ is calculated from held-out residuals using distance-band weights, with the threshold set by a semivariogram range \ref{app}
\begin{equation}
I=\frac{N}{W}
\frac{\sum_i\sum_j w_{ij}(e_i-\bar e)(e_j-\bar e)}
{\sum_i(e_i-\bar e)^2}.
\label{eq:morans-i}
\end{equation}
It is used after fitting, not for feature selection. The same rule is applied to RQ2a error
reduction, RQ2b residuals and the Elastic Net, XGBoost and GNNWR residuals in RQ3.

Neural training variation is assessed over five seeds for the selected RQ1 models. XGBoost and the
DNN are also compared fold by fold on identical test folds. The physics-weight ablation tests
weights 0, 1 and 2. For PINN-$k$, parameter separability is assessed by the pooled correlation
between learned $y_{\max}$ and $k$. Its confidence interval uses the same 2,000 compartment
bootstrap resamples. These checks distinguish predictive accuracy from parameter stability.

\section{Reproducibility}

Split assignments, feature lists and model settings are saved with each run. Seed 42 is used unless
a stated reseed is part of the analysis. Full feature membership, hyperparameter searches,
secondary split results and detailed outlier tables are retained in the Appendix. Methodological
limitations and their effect on interpretation are evaluated after the results, before the final
conclusions.
